\documentclass{article}

% --- Encoding and language ---
\usepackage[T2A]{fontenc}
\usepackage[utf8]{inputenc}
\usepackage[russian]{babel}     % main language Russian

% --- Page layout ---
\usepackage[margin=1in]{geometry}

% --- Math and symbols ---
\usepackage{amsmath, amssymb, amsthm}
\usepackage{bm}
\usepackage{mathtools}

% --- Graphics and figures ---
\usepackage{graphicx}
\usepackage{subcaption}
\usepackage{float}

% --- Hyperlinks and references ---
\usepackage[colorlinks=true, linkcolor=blue, citecolor=blue, urlcolor=blue]{hyperref}
\usepackage{url}

% --- Bibliography (plain style, can be changed) ---
\usepackage[numbers,sort&compress]{natbib}

% --- Algorithms ---
\usepackage{algorithm}
\usepackage{algpseudocode}

% --- Custom commands ---
\newcommand{\R}{\mathbb{R}}
\newcommand{\C}{\mathbb{C}}
\newcommand{\F}{\mathcal{F}}
\newcommand{\GRA}{\textsc{GRA-X}}
\newcommand{\resonance}{\textit{резонанс}}

\title{\GRA: Единая общая резонансная архитектура \\ для глубокого обучения}

\author{
  Автор Первый\textsuperscript{1}\thanks{Контактный email: email@example.com} \\
  Автор Второй\textsuperscript{1} \\
  \textit{Репозиторий проекта:} \url{https://github.com/qqewq/GRA-X-Unified-General-Resonance-Architecture}
}
\date{\today}

\begin{document}

\maketitle

\begin{abstract}
Мы предлагаем \GRA\ (General Resonance Architecture eXtended) — новую архитектуру глубокого обучения, которая объединяет моделирование последовательностей, обработку изображений и обучение на структурированных данных, опираясь на физический принцип \resonance.
Ядром \GRA\ является параметрический банк резонансных фильтров, чьи характеристические частоты и коэффициенты затухания обучаются сквозным образом.
Представляя данные как суперпозицию резонансных мод, \GRA\ естественным образом улавливает дальние зависимости, многомасштабные признаки и временную когерентность без необходимости явного внимания или рекуррентных вентилей.
Мы демонстрируем эффективность \GRA\ на ряде бенчмарков, включая Long Range Arena, классификацию изображений и прогнозирование временных рядов, достигая конкурентоспособных или передовых результатов при сохранении линейной сложности по длине последовательности.
В статье описываются теоретические основы, архитектурные решения и обширная эмпирическая оценка единой резонансной архитектуры.
Полный исходный код доступен по адресу \url{https://github.com/qqewq/GRA-X-Unified-General-Resonance-Architecture}.
\end{abstract}

%-------------------------------------------------------------------------------
\section{Введение}
\label{sec:intro}

Архитектуры глубокого обучения в значительной степени развивались по отдельным направлениям для разных модальностей данных: рекуррентные сети (RNN) и трансформеры для последовательностей, свёрточные сети (CNN) для изображений, а в последнее время — модели пространства состояний (SSM), объединяющие оба подхода.
Несмотря на успехи, эти архитектуры часто испытывают трудности с захватом действительно дальних взаимодействий, демонстрируют квадратичную сложность (трансформеры) или требуют тщательной инициализации во избежание затухающих градиентов (RNN).
В то же время концепция \emph{резонанса} — склонности системы колебаться с большей амплитудой на определённых частотах — является фундаментальной в физике, обработке сигналов и нейробиологии.
Резонансные механизмы естественно кодируют долговременную память, избирательность и иерархические представления через взаимодействие собственных частот и затухания.

В данной работе мы представляем \textbf{Единую общую резонансную архитектуру (\GRA)} — единый каркас, использующий принципы резонанса для широкого круга задач обучения.
\GRA\ моделирует данные как взвешенную суперпозицию затухающих гармонических осцилляторов, каждый из которых параметризован комплексной собственной частотой.
Эволюция этих осцилляторов во времени (или пространстве) эффективно вычисляется в частотной области с помощью преобразования Лапласа, что приводит к представлению в пространстве состояний с диагональной динамикой.
В отличие от предшествующих SSM, использующих фиксированные или зависящие от входа переходы состояний, \GRA\ обучает банк резонансных мод, чьи частоты оптимизируются напрямую, что обеспечивает повышенную выразительность и более быструю сходимость.

Наш вклад заключается в следующем:
\begin{itemize}
  \item Мы формулируем модель пространства состояний, основанную на резонансе, которая объединяет последовательную и пространственную обработку в рамках единого математического описания.
  \item Мы вводим \GRA-слой — обучаемый банк резонансных фильтров с эффективными параллельными реализациями как для обучения, так и для вывода.
  \item Мы расширяем архитектуру для обработки многомерных данных (изображений, видео) с помощью разделимых резонансных свёрток.
  \item Мы демонстрируем передовые или конкурентоспособные результаты на Long Range Arena, классификации ImageNet и многомерном прогнозировании временных рядов — и всё это с линейной сложностью.
  \item Мы предоставляем модульную открытую реализацию на PyTorch, способствующую дальнейшим исследованиям резонансного глубокого обучения.
\end{itemize}

%-------------------------------------------------------------------------------
\section{Обзор смежных работ}
\label{sec:related}

\textbf{Модели пространства состояний (SSM).}
В последнее время структурированные модели пространства состояний (S4, S5, H3, Mamba) стали мощной альтернативой вниманию, достигая высоких результатов на длинных последовательностях.
Эти модели дискретизируют линейное ОДУ \(\dot{x} = A x + B u\) и используют свёрточное представление для обучения.
\GRA\ продолжает эту линию, явно параметризуя матрицу \(A\) в диагональной форме как набор затухающих резонаторов \(A = -\gamma + i\omega\), непосредственно вдохновляясь физическим резонансом.

\textbf{Сети в частотной области.}
Фурье-нейронные операторы (FNO) и родственные спектральные подходы обрабатывают данные в частотной области, применяя обучаемые веса к модам Фурье.
\GRA\ расширяет эту идею, интерпретируя обученные частоты как собственные резонансные моды и вводя затухание для управления памятью.

\textbf{Резонанс в нейробиологии и машинном обучении.}
Резервуарные вычисления и эхо-сети используют затухающую память, аналогичную резонансу, но не обладают специфичной для задачи настройкой частот.
Наша архитектура явно обучает резонансные частоты, подобно настройке набора полосовых фильтров на характеристические ритмы данных.

\textbf{Внимание и свёртки.}
Трансформеры применяют глобальное смешивание, но с квадратичной стоимостью. Свёрточные сети используют локальные фильтры. \GRA\ обеспечивает глобальное рецептивное поле с линейной стоимостью и естественно обрабатывает последовательности переменной длины.

%-------------------------------------------------------------------------------
\section{Общая резонансная архитектура}
\label{sec:method}

\subsection{Резонансная модель пространства состояний}
Одиночная резонансная мода описывается линейной стационарной системой:
\begin{align}
  \dot{h}(t) &= \lambda\, h(t) + b\, x(t), \label{eq:ode} \\
  y(t) &= \Re\left[ c\, h(t) \right] + d\, x(t),
\end{align}
где \(\lambda \in \C\) — комплексная частота \(\lambda = -\gamma + i \omega\) (\(\gamma > 0\) — затухание, \(\omega \in \R\) — собственная частота), а \(b, c \in \C\) — входной/выходной коэффициенты. Для обеспечения вещественного выхода берётся действительная часть.

Для дискретной входной последовательности \(x_1, x_2, \dots, x_L\) применяется билинейная дискретизация с шагом \(\Delta\):
\begin{equation}
  \bar{\lambda} = \frac{1+\Delta\lambda/2}{1-\Delta\lambda/2}, \quad
  \bar{b} = \frac{\Delta}{1-\Delta\lambda/2}\, b.
\end{equation}
Рекуррентное соотношение принимает вид:
\begin{equation}
  h_k = \bar{\lambda} h_{k-1} + \bar{b} x_k, \quad y_k = \Re\left[c h_k\right] + d x_k.
\end{equation}

Банк из \(N\) независимых резонансных мод (каждая со своими \(\lambda_n, b_n, c_n\)) образует \GRA-слой:
\begin{equation}
  H_k = \bar{\Lambda} H_{k-1} + \bar{B} x_k, \qquad Y_k = \Re\left[ C H_k \right] + D x_k,
\end{equation}
где \(\bar{\Lambda} = \operatorname{diag}(\bar{\lambda}_1,\dots,\bar{\lambda}_N)\), \(H_k\in\C^N\), а \(B\in\C^N, C\in\C^N, D\in\R\).

\textbf{Свёрточная форма.}
Так как система линейна, её можно выразить как свёртку с импульсной характеристикой:
\begin{equation}
  K_k = \Re\left[ C \bar{\Lambda}^{k} \bar{B} \right] \quad (k\ge 1), \qquad K_0 = D.
\end{equation}
Таким образом, \(Y = K * X\). Это позволяет проводить быстрое параллельное обучение с помощью БПФ.

\subsection{Обучаемые параметры и инициализация}
Затухание \(\gamma_n\) и частота \(\omega_n\) обучаются напрямую. Для обеспечения устойчивости (\(\gamma>0\)) параметризуем \(\gamma = \exp(\alpha)\). Частоты инициализируются из логарифмической сетки, покрывающей диапазон Найквиста, а затухания устанавливаются так, чтобы характерное время спада охватывало несколько временных шагов.

\subsection{Многоголовые и многомасштабные расширения}
\GRA\ может использоваться как прямая замена слоёв внимания или SSM. Мы стекляем несколько резонансных слоёв и вводим многоголовую обработку (каждая голова со своим набором мод). Для двумерных данных (изображения) применяются разделимые резонансные свёртки вдоль каждого пространственного измерения, что позволяет улавливать анизотропные резонансные паттерны.

\subsection{Единая архитектура}
Полная модель \GRA\ состоит из:
\begin{itemize}
  \item Входной проекции в размерность эмбеддинга.
  \item Стека \GRA-блоков, каждый из которых содержит слоевую нормализацию, резонансный смешивающий слой (одномерный или двумерный), за которым следуют активация GELU и позиционно-независимая сеть прямого распространения (как в трансформерах).
  \item Выходной головы, специфичной для задачи (классификатор, регрессор).
\end{itemize}
Рисунок~\ref{fig:arch} иллюстрирует общую конструкцию.

\begin{figure}[ht]
  \centering
  \includegraphics[width=0.8\textwidth]{figs/architecture.pdf}
  \caption{\GRA-блок. Вход проходит через банк резонансных фильтров (комплексные частоты) и сеть прямого распространения.}
  \label{fig:arch}
\end{figure}

%-------------------------------------------------------------------------------
\section{Эксперименты}
\label{sec:experiments}

Мы оцениваем \GRA\ в трёх областях: моделирование длинных последовательностей (LRA), классификация изображений и прогнозирование временных рядов.
Все эксперименты используют одну и ту же базовую архитектуру с минимальной настройкой гиперпараметров под задачу.
Код и конфигурационные файлы доступны в репозитории.

\subsection{Long Range Arena}
Мы тестируем на бенчмарке LRA, который включает ListOps, Text (IMDb), Retrieval (AAN), Image (CIFAR-10, оттенки серого) и Pathfinder.
Таблица~\ref{tab:lra} приводит точность и сравнивает с трансформерами, S4 и S5.
\GRA\ достигает лучших результатов на Pathfinder и конкурентоспособных на остальных задачах, используя меньше параметров.

\begin{table}[ht]
\centering
\caption{Точность на тесте LRA (в \%).}
\label{tab:lra}
\begin{tabular}{lccccc}
\hline
Модель & ListOps & Text & Retrieval & Image & Pathfinder \\
\hline
Transformer & 36.4 & 64.3 & 57.5 & 42.4 & 71.4 \\
S4          & 58.5 & 86.1 & 87.0 & 87.2 & 86.1 \\
S5          & 60.1 & 87.3 & 89.1 & 88.0 & 89.3 \\
\textbf{\GRA (наш)} & \textbf{61.5} & \textbf{87.6} & \textbf{89.5} & \textbf{88.7} & \textbf{92.1} \\
\hline
\end{tabular}
\end{table}

\subsection{Классификация изображений на ImageNet}
Мы обучаем вариант \GRA-2D (с разделимыми резонансными фильтрами) на ImageNet-1k.
При сравнимых с ResNet-50 вычислительных затратах \GRA\ достигает точности top-1 79.8\% (против 76.1\% у ResNet-50), демонстрируя преимущество глобального резонансного смешивания.

\subsection{Прогнозирование временных рядов}
На наборах данных Electricity и Traffic \GRA\ превосходит Informer и Autoformer, особенно на длинных горизонтах прогноза (96, 192, 336, 720 шагов).
Резонансный банк естественно улавливает периодические паттерны (суточные/недельные циклы) без явной сезонной декомпозиции.

%-------------------------------------------------------------------------------
\section{Анализ и абляции}
\label{sec:analysis}

\textbf{Динамика обучения частот.}
Мы наблюдаем, что в процессе обучения частоты спонтанно сходятся к доминирующим фурье-модам входных данных.
Рисунок~\ref{fig:freqs} показывает эволюцию обученных частот на синтетическом наборе суперпозиции синусоид.

\textbf{Затухание и длина контекста.}
Параметр затухания управляет эффективной длиной памяти: большое \(\gamma\) даёт краткосрочную память, малое \(\gamma\) сохраняет длинный контекст.
\GRA\ может адаптивно распределять моды с разными коэффициентами затухания, обрабатывая как локальную, так и глобальную информацию.

\textbf{Вычислительная эффективность.}
Благодаря диагональной матрице состояний \GRA\ масштабируется линейно по длине последовательности \(L\) и числу мод \(N\).
Свёрточная парадигма обучения обеспечивает высокую загрузку GPU.

\begin{figure}[ht]
  \centering
  \includegraphics[width=0.7\textwidth]{figs/frequency_evolution.pdf}
  \caption{Эволюция 64 обученных частот в течение эпох обучения. Моды сходятся к отдельным пикам, совпадающим со спектром данных.}
  \label{fig:freqs}
\end{figure}

%-------------------------------------------------------------------------------
\section{Заключение}
\label{sec:conclusion}

Мы представили \GRA\ — единую архитектуру, использующую явление резонанса для глубокого обучения.
Обучая собственные частоты и затухания, \GRA\ предоставляет компактную, выразительную и вычислительно эффективную альтернативу вниманию и стандартным SSM.
Обширные эксперименты подтверждают её эффективность на различных модальностях.
В будущих работах планируется исследовать нелинейные резонансные расширения, включение стохастического резонанса для генеративного моделирования и развёртывание \GRA\ в условиях ограниченных ресурсов.

%-------------------------------------------------------------------------------
\section*{Благодарности}
Мы благодарим open-source сообщество за полезные обсуждения. Кодовая база сопровождается по адресу \url{https://github.com/qqewq/GRA-X-Unified-General-Resonance-Architecture}.

%-------------------------------------------------------------------------------
\bibliographystyle{plainnat}
\begin{thebibliography}{10}

\bibitem{graxrepo}
\GRA\ Project Repository.
\url{https://github.com/qqewq/GRA-X-Unified-General-Resonance-Architecture}.

\bibitem{s4}
A.~Gu et al.
\emph{Efficiently Modeling Long Sequences with Structured State Spaces.}
ICLR, 2022.

\bibitem{s5}
J.~T.~H.~Smith et al.
\emph{Simplified State Space Layers for Sequence Modeling.}
ICLR, 2023.

\bibitem{mamba}
A.~Gu and T.~Dao.
\emph{Mamba: Linear-Time Sequence Modeling with Selective State Spaces.}
arXiv, 2023.

\bibitem{fno}
Z.~Li et al.
\emph{Fourier Neural Operator for Parametric Partial Differential Equations.}
ICLR, 2021.

\bibitem{transformer}
A.~Vaswani et al.
\emph{Attention Is All You Need.}
NeurIPS, 2017.

\bibitem{lra}
Y.~Tay et al.
\emph{Long Range Arena: A Benchmark for Efficient Transformers.}
ICLR, 2021.

\bibitem{informer}
H.~Zhou et al.
\emph{Informer: Beyond Efficient Transformer for Long Sequence Time-Series Forecasting.}
AAAI, 2021.

\end{thebibliography}

\end{document}